\chapter{Lezione 14 - 7/11}

\section{Applicazione Lineare Associata ad un Endomorfismo}

\Teorema{Applicazione Lineare Associata ad un Endomorfismo}{

    Sia \(T: V \to V\) un endomorfismo (applicazione lineare da \(V\) in se stesso), con \(\dim V = n\).
    Siano \(\mathcal{B}\) e \(\bar{\mathcal{B}}\) due basi diverse di \(V\).
    
    Possiamo associare a \(T\) due matrici diverse, una per ogni base:
    \begin{itemize}
        \item \(A = M_{\mathcal{B}}(T)\) (Matrice associata rispetto alla base \(\mathcal{B}\))
        \item \(\bar{A} = M_{\bar{\mathcal{B}}}(T)\) (Matrice associata rispetto alla base \(\bar{\mathcal{B}}\))
    \end{itemize}
    
   
    \Dimostrazione{
        \textbf{Dimostrazione della relazione (Change of Basis):}
    
    Sia \(P = M_{\mathcal{B}, \bar{\mathcal{B}}}(id_V)\) la matrice di passaggio dalla base \(\mathcal{B}\) alla base \(\bar{\mathcal{B}}\).
    Questa matrice trasforma le coordinate "vecchie" (\(X\)) in coordinate "nuove" (\(\bar{X}\)):
    \[ \bar{X} = P X \quad \text{e analogamente} \quad \bar{Y} = P Y \]
    
    Preso un vettore \(u \in V\), consideriamo le sue coordinate e quelle della sua immagine \(T(u)\):
    \begin{itemize}
        \item In base \(\mathcal{B}\): \(\phi_{\mathcal{B}}(u) = X\) e \(\phi_{\mathcal{B}}(T(u)) = Y\). Vale la relazione:
        \[ Y = A X \]
        
        \item In base \(\bar{\mathcal{B}}\): \(\phi_{\bar{\mathcal{B}}}(u) = \bar{X}\) e \(\phi_{\bar{\mathcal{B}}}(T(u)) = \bar{Y}\). Vale la relazione:
        \[ \bar{Y} = \bar{A} \bar{X} \]
    \end{itemize}
    
    Sostituiamo le relazioni di passaggio (\(\bar{X}=PX, \bar{Y}=PY\)) nell'equazione della base \(\bar{\mathcal{B}}\):
    \[
    \bar{Y} = \bar{A} \bar{X} \implies P Y = \bar{A} (P X)
    \]
    Moltiplichiamo a sinistra per \(P^{-1}\) per isolare \(Y\):
    \[
    Y = P^{-1} \bar{A} P X
    \]
    Confrontando questa equazione con \(Y = A X\), per l'unicità della matrice associata, otteniamo:
    \[
    A = P^{-1} \bar{A} P
    \]
    
    \textbf{Conclusione:} Due matrici che rappresentano lo stesso endomorfismo rispetto a basi diverse sono \textbf{simili}.
    }
    
}

\section{Matrici Simili}
\Definizione{Matrici Simili}{
    Due matrici quadrate \(A, \bar{A} \in M_n(K)\) si dicono \textbf{simili} se esiste una matrice invertibile \(P \in M_n(K)\) tale che:
    \[
    A = P^{-1} \bar{A} P \quad (\text{oppure } \bar{A} = P A P^{-1})
    \]
}
\section{Autovalori, Autovettori e Autospazi}

\Definizione{Autovalore}{
    Sia \(T: V \to V\) un endomorfismo su un campo \(K\).
    Uno scalare \(\lambda \in K\) si dice \textbf{autovalore} di \(T\) \(\iff\) l'insieme:
    \[ U_\lambda = \{u \in V \mid T(u) = \lambda u\} \]
    contiene vettori non nulli, ovvero \(U_\lambda \neq \{\bar{0}\}\).

    \textbf{Osservazione:}
    Consideriamo il caso \(\lambda = 0\).
    \[ U_0 = \{u \in V \mid T(u) = 0 \cdot u = \bar{0}\} = \ker(T) \]
    Quindi: \(0\) è autovalore di \(T\) \(\iff \ker(T) \neq \{\bar{0}\} \iff T\) non è iniettiva.
}

\Proposizione{L'Autospazio è un Sottospazio}{
    Per ogni \(\lambda \in K\), l'insieme \(U_\lambda\) è un sottospazio vettoriale di \(V\).
    
    \Dimostrazione{
        Verifichiamo le tre proprietà dei sottospazi:
        \begin{itemize}
            \item \textbf{Contiene il vettore nullo:}
            Poiché \(T\) è lineare, \(T(\bar{0}) = \bar{0}\). Inoltre \(\lambda \cdot \bar{0} = \bar{0}\).
            Quindi \(T(\bar{0}) = \lambda \bar{0} \implies \bar{0} \in U_\lambda\).

            \item \textbf{Chiusura rispetto alla somma:}
            Siano \(u, v \in U_\lambda\). Ciò significa che \(T(u) = \lambda u\) e \(T(v) = \lambda v\).
            Calcoliamo \(T(u+v)\):
            \[ T(u+v) = T(u) + T(v) = \lambda u + \lambda v = \lambda(u+v) \]
            Quindi \((u+v) \in U_\lambda\).

            \item \textbf{Chiusura rispetto al prodotto per scalare:}
            Siano \(u \in U_\lambda\) e \(\alpha \in K\).
            Calcoliamo \(T(\alpha u)\):
            \[ T(\alpha u) = \alpha T(u) = \alpha (\lambda u) = (\alpha \lambda) u = (\lambda \alpha) u = \lambda (\alpha u) \]
            Quindi \(\alpha u \in U_\lambda\).
        \end{itemize}
    }
}

\Definizione{Autospazio e Autovettore}{
    Se \(\lambda\) è un autovalore di \(T\), allora il sottospazio \(U_\lambda\) si definisce \textbf{autospazio} relativo a \(\lambda\).
    
    I vettori \textbf{non nulli} appartenenti a \(U_\lambda\) si dicono \textbf{autovettori} relativi all'autovalore \(\lambda\).
}


\BoxBlue{Osservazione}{Calcolo di Autovalori e Autospazi (in \(V\) f.g.)}{
    Sia \(T: V \to V\) un endomorfismo e sia \(\mathcal{B}\) una base di \(V\).
    Sia \(A = M_{\mathcal{B}}(T)\) la matrice associata.
    Poniamo \(X = \phi_{\mathcal{B}}(u)\) (il vettore colonna delle coordinate).

    \textbf{1. Ricerca degli Autovalori (Polinomio Caratteristico)}
    
    \(\lambda \in K\) è autovalore di \(T\)
    \[ \iff \exists u \in V \setminus \{\bar{0}\} : T(u) = \lambda u \]
    Passando alle coordinate (grazie all'isomorfismo \(\phi_{\mathcal{B}}\)):
    \[ \iff \exists X \in K^n \setminus \{\bar{0}\} : A X = \lambda X \]
    Portando tutto a sinistra e raccogliendo la \(X\) (usando la matrice identità \(I_n\)):
    \[ \iff A X - \lambda I_n X = \bar{0} \]
    \[ \iff \exists X \neq \bar{0} : (A - \lambda I_n) X = \bar{0} \]
    
    Un sistema omogeneo ammette soluzioni non nulle (\(X \neq \bar{0}\)) se e solo se la matrice dei coefficienti \textbf{non è invertibile}, ovvero ha determinante nullo:
    \[ \iff \det(A - \lambda I_n) = 0 \]
    
    Questa equazione è detta \textbf{Equazione Caratteristica}.
    Il polinomio \(p_A(\lambda) = \det(A - \lambda I_n)\) è detto \textbf{Polinomio Caratteristico}.

    \bigskip
    \textbf{2. Ricerca degli Autospazi}
    
    Per ogni \(\bar{\lambda}\) autovalore trovato (radice del polinomio), l'autospazio \(U_{\bar{\lambda}}\) si trova risolvendo il sistema lineare omogeneo associato:
    \[ \Sigma_0: (A - \bar{\lambda}I_n)X = \bar{0} \]
    
    L'insieme delle soluzioni \(S_0\) di questo sistema rappresenta le coordinate dei vettori dell'autospazio:
    \[ S_0 = \phi_{\mathcal{B}}(U_{\bar{\lambda}}) \]
    (Per trovare i vettori "veri" in \(V\), basta applicare l'isomorfismo inverso ai vettori soluzione \(X\)).
}

\Proposizione{Invarianza del Polinomio Caratteristico per Similitudine}{
    Siano \(A, \bar{A} \in M_n(K)\).
    Se \(A\) e \(\bar{A}\) sono simili (cioè esiste \(P\) invertibile tale che \(A = P^{-1} \bar{A} P\)), allora hanno lo stesso polinomio caratteristico:
    \[
    |A - \lambda I_n| = |\bar{A} - \lambda I_n|
    \]
    (Di conseguenza, hanno gli stessi autovalori con la stessa molteplicità algebrica).

    \Dimostrazione{
        Partiamo dal polinomio caratteristico di \(A\):
        \[ |A - \lambda I_n| \]
        Sostituiamo \(A\) con la sua espressione simile \(P^{-1} \bar{A} P\):
        \[ = |P^{-1} \bar{A} P - \lambda I_n| \]
        
        Qui usiamo il trucco fondamentale: scriviamo la matrice identità come \(I_n = P^{-1} I_n P\). 
        L'espressione diventa:
        \[ = |P^{-1} \bar{A} P - \lambda (P^{-1} I_n P)| \]
        
        Poiché \(\lambda\) è uno scalare, può essere spostato all'interno del prodotto matriciale:
        \[ = |P^{-1} \bar{A} P - P^{-1} (\lambda I_n) P| \]
        
        Ora possiamo raccogliere \(P^{-1}\) a sinistra e \(P\) a destra (proprietà distributiva delle matrici):
        \[ = |P^{-1} (\bar{A} - \lambda I_n) P| \]
        
        Applichiamo il \textbf{Teorema di Binet} (il determinante del prodotto è il prodotto dei determinanti):
        \[ = |P^{-1}| \cdot |\bar{A} - \lambda I_n| \cdot |P| \]
        
        Poiché i determinanti sono scalari (numeri in \(K\)), il prodotto è commutativo. Possiamo riordinare i termini:
        \[ = |P^{-1}| \cdot |P| \cdot |\bar{A} - \lambda I_n| \]
        
        Sapendo che \(|P^{-1}| = |P|^{-1} = \frac{1}{|P|}\), il prodotto dei determinanti di \(P\) e \(P^{-1}\) fa 1:
        \[ = 1 \cdot |\bar{A} - \lambda I_n| \]
        \[ = |\bar{A} - \lambda I_n| \]
        
        La tesi è dimostrata.
    }
}

\bigskip

\Proposizione{Indipendenza degli Autovettori relativi ad autovalori distinti}{
    Sia \(T\) un endomorfismo su uno spazio vettoriale \(V\) sul campo \(K\).
    Siano \(\lambda_1, \dots, \lambda_h\) autovalori di \(T\) \textbf{a due a due distinti} (\(\lambda_i \neq \lambda_j\) per \(i \neq j\)).
    
    Per ogni autovalore \(\lambda_i\), consideriamo un autovettore non nullo \(u_i \in U_{\lambda_i} \setminus \{\bar{0}\}\).
    
    Allora l'insieme \(\{u_1, \dots, u_h\}\) è \textbf{linearmente indipendente}.

    \Dimostrazione{
        \begin{itemize}
            \item \textbf{Preliminare: Intersezione nulla.}
            Innanzitutto osserviamo che \(U_{\lambda_i} \cap U_{\lambda_j} = \{\bar{0}\}\) con \(i \neq j\).
            Infatti, sia \(u \in U_{\lambda_i} \cap U_{\lambda_j}\).
            Allora \(T(u) = \lambda_i u\) e \(T(u) = \lambda_j u\).
            Uguagliando le due espressioni:
            \[ \lambda_i u = \lambda_j u \implies (\lambda_i - \lambda_j)u = \bar{0} \]
            Poiché gli autovalori sono distinti, \((\lambda_i - \lambda_j) \neq 0\), quindi deve essere necessariamente \(u = \bar{0}\).

            \item \textbf{Dimostrazione per induzione su \(h\).}
            
            \textbf{Base dell'induzione (\(h=1\)):}
            L'insieme \(\{u_1\}\) contiene un solo vettore non nullo, quindi è linearmente indipendente.
            
            \textbf{Passo induttivo (\(h-1 \implies h\)):}
            Supponiamo per ipotesi induttiva che \(h-1\) autovettori relativi ad autovalori distinti siano L.I.
            Consideriamo una combinazione lineare nulla di \(h\) autovettori:
            \[
            \alpha_1 u_1 + \dots + \alpha_{h-1} u_{h-1} + \alpha_h u_h = \bar{0} \quad (*).
            \]
            Dobbiamo dimostrare che tutti gli scalari \(\alpha_i\) sono nulli.
            
            \textbf{Passo A:} Moltiplichiamo l'equazione \((*)\) per lo scalare \(\lambda_h\):
            \[
            \lambda_h \alpha_1 u_1 + \dots + \lambda_h \alpha_{h-1} u_{h-1} + \lambda_h \alpha_h u_h = \bar{0} \quad (1)
            \]
            
            \textbf{Passo B:} Applichiamo l'applicazione lineare \(T\) all'equazione \((*)\):
            \[
            T(\alpha_1 u_1 + \dots + \alpha_h u_h) = T(\bar{0}) = \bar{0}
            \]
            Per la linearità di \(T\) e la definizione di autovettore (\(T(u_i) = \lambda_i u_i\)):
            \[
            \alpha_1 \lambda_1 u_1 + \dots + \alpha_{h-1} \lambda_{h-1} u_{h-1} + \alpha_h \lambda_h u_h = \bar{0} \quad (2)
            \]
            
            \textbf{Passo C:} Sottraiamo l'equazione (1) dall'equazione (2):
            \[
            \alpha_1 (\lambda_1 - \lambda_h) u_1 + \dots + \alpha_{h-1} (\lambda_{h-1} - \lambda_h) u_{h-1} + \underbrace{\alpha_h (\lambda_h - \lambda_h) u_h}_{=0} = \bar{0}
            \]
            Otteniamo una combinazione lineare di soli \(h-1\) vettori:
            \[
            \sum_{i=1}^{h-1} \alpha_i (\lambda_i - \lambda_h) u_i = \bar{0}
            \]
            Per l'\textbf{ipotesi induttiva}, i vettori \(\{u_1, \dots, u_{h-1}\}\) sono linearmente indipendenti. Quindi, tutti i coefficienti di questa combinazione devono essere nulli:
            \[
            \alpha_i (\lambda_i - \lambda_h) = 0 \quad \forall i = 1, \dots, h-1.
            \]
            Poiché gli autovalori sono distinti, \(\lambda_i \neq \lambda_h\) per \(i < h\), quindi il termine tra parentesi non è mai zero. Ne consegue che:
            \[
            \alpha_1 = 0, \ \alpha_2 = 0, \ \dots, \ \alpha_{h-1} = 0.
            \]
            
            Sostituendo questi valori nell'equazione originale \((*)\), rimane solo \(\alpha_h u_h = \bar{0}\). Dato che \(u_h \neq \bar{0}\), deve essere \(\alpha_h = 0\).
            Tutti i coefficienti sono nulli, quindi l'insieme è L.I.
        \end{itemize}
    }
}

\BoxBlue{Corollario}{Somma Diretta degli Autospazi}{
    Sia \(T: V \to V\) un endomorfismo.
    Siano \(\lambda_1, \dots, \lambda_h\) autovalori a due a due distinti e siano \(U_{\lambda_1}, \dots, U_{\lambda_h}\) i relativi autospazi.
    
    Allora la somma degli autospazi è \textbf{diretta}:
    \[ U_{\lambda_1} + \dots + U_{\lambda_h} = U_{\lambda_1} \oplus \dots \oplus U_{\lambda_h} \]
    
    Questo significa che per ogni \(j \in \{1, \dots, h\}\), l'intersezione dell'autospazio \(j\)-esimo con la somma degli altri è banale:
    \[
    U_{\lambda_j} \cap \left( \sum_{i \neq j} U_{\lambda_i} \right) = \{\bar{0}\}
    \]

    \Dimostrazione{
        Senza perdita di generalità, dimostriamolo per \(j=1\).
        Dobbiamo dimostrare che \(U_{\lambda_1} \cap (U_{\lambda_2} + \dots + U_{\lambda_h}) = \{\bar{0}\}\).
        
        Sia \(u\) un vettore appartenente a questa intersezione:
        \[ u \in U_{\lambda_1} \quad \text{e} \quad u \in (U_{\lambda_2} + \dots + U_{\lambda_h}) \]
        
        Poiché \(u\) appartiene alla somma, esistono dei vettori \(u_2 \in U_{\lambda_2}, \dots, u_h \in U_{\lambda_h}\) tali che:
        \[ u = u_2 + \dots + u_h \]
        Portando \(u\) a destra, possiamo scrivere una combinazione lineare nulla:
        \[ \bar{0} = -u + u_2 + \dots + u_h \]
        
        Ora ragioniamo per assurdo.
        Supponiamo che \(u \neq \bar{0}\) (o che qualche \(u_i \neq \bar{0}\)).
        \begin{itemize}
            \item \(u\) è un autovettore relativo a \(\lambda_1\) (o è nullo).
            \item Ogni \(u_i\) è un autovettore relativo a \(\lambda_i\) (o è nullo).
        \end{itemize}
        L'equazione \((-u) + u_2 + \dots + u_h = \bar{0}\) rappresenterebbe una combinazione lineare non banale di autovettori relativi ad autovalori \textbf{distinti} che dà il vettore nullo.
        
        Ma per la \textbf{Proposizione precedente}, autovettori relativi ad autovalori distinti sono \textbf{linearmente indipendenti}. L'unica combinazione lineare che dà zero è quella con tutti i vettori nulli.
        
        Quindi deve essere necessariamente:
        \[ u = \bar{0}, \quad u_2 = \bar{0}, \quad \dots, \quad u_h = \bar{0} \]
        
        Poiché l'unico vettore nell'intersezione è il vettore nullo, la somma è diretta.
    }
}